\hypertarget{appendix-f82e270791c8}{}
\section*{Accumulated rounding error}
\section{Accumulated rounding error}\label{accumulated-rounding-error}

\textbf{Curriculum depth:} Introduced · \textbf{First seen:} C005 · \textbf{Latest development:} Phase 10

\subsection{Definition and mechanism}\label{definition-and-mechanism}

Under suitable nearest normal-range assumptions, products of rounding factors are bounded using gamma\_n=nu/(1−nu), nu\textless1. A sequential n-term sum has absolute error bounded by gamma\_(n−1) times sum of input magnitudes. Cancellation can make its relative error much larger.

\subsection{Prerequisites}\label{prerequisites}

\begin{itemize}
\tightlist
\item
  \hyperref[appendix-bc44f7b42f42]{Unit roundoff} --- The local normal-rounding bound supplies the error-analysis scale.
\item
  \hyperref[appendix-3ae74e89b4e6]{Absolute and relative error} --- Discrepancy, magnitude and reference units define sensitivity and threshold significance.
\end{itemize}

\subsection{Used by}\label{used-by}

\begin{itemize}
\item
  \hyperref[appendix-df318da57389]{Compensated and pairwise summation} --- uses this mechanism to establish its numerical contract.
\item
  \href{}{Online moments and Welford accumulation (02-Foundations-II, appendix-33237f0f618b)}
\item
  \href{}{Discretization, truncation, and roundoff (02-Foundations-II, appendix-6b0648b2b8e8)} \#\# Recall and next depth
\end{itemize}

Reconstruct the definition's example and name the assumptions that make it valid.

\textbf{Pending:} Apply bounds to actual dependency paths, not a global count of source operators; subnormals/range failures require other analysis.

\subsection{Source and navigation}\label{source-and-navigation}

\href{Continuous\%20Notes/Computing\%20Foundations\%20-\%20From\%20Bits\%20to\%20C++\%20and\%20Systems\%20APIs\#When\%20arithmetic\%20meets\%20uncertainty}{M005 teaching chapter} · \hyperref[appendix-fd090d982a15]{Part definitions} · \href{}{Dependency Map (02-Foundations-II, appendix-4ad47bdd269a)}

\subsection{Full teaching explanation}\label{full-teaching-explanation}

\begin{itemize}
\tightlist
\item
  \href{Continuous\%20Notes/Computing\%20Foundations\%20-\%20From\%20Bits\%20to\%20C++\%20and\%20Systems\%20APIs\#9.\%20Finite\%20samples\%20need\%20stable\%20descriptive\%20statistics}{Phase 10 teaching section}
\item
  \hyperref[appendix-cb5af353ef77]{Visual Recall/10 - Probability and Monte Carlo - Visual Recall.md} · Phase 10 visual recall
\end{itemize}
